% page 359 of the facsimile (image p-358 of the scan)
% block 162.6 x 237.7 mm ; left margin 39.4 mm
\pgc{17.780mm}{0.000mm}{
\l{\cel{48}351}
\l{}
\l{\sou{makrosporangio}. (\sou{biol.})\cel{2}Sporangio qua produktas la makro-}
\l{\cel{3}sporo. -}
\l{}
\l{\sou{makulo}. Marko qua sordidigis ulo.- II. Marko naturala sur la}
\l{\cel{3}pelo di la animali. - III. Altereso, en un loko, di la koloro}
\l{\cel{3}di la ensemblo. - (\sou{anke metaf.})\cel{2}DEFIS.}
\l{}
\l{\sou{makulaturo}. (\sou{imprim-arto}). Papero-folio qua uzesis por recevar}
\l{\cel{3}la quanto ecesanta di la imprim-inko - o : imprimo-folio en}
\l{\cel{3}qua la tipi esas desordinita pro mala imprimo. - DEFIRS.}
\l{}
\l{\sou{mala.}\cel{2}I. Ye maniero regretinda, penigiva. - II. Ye maniero}
\l{\cel{3}neperfekta, defektoza. - III. Ye maniero kontrea a la lego}
\l{\cel{3}etikala, morala. - IV. Kontrea a lo devata. - EFIS.}
\l{}
\l{\sou{malada}. I.Qua subisis kelka altero di sua organi, kelka par-}
\l{\cel{3}turbo di sua funcioni organismala. - II. (\sou{metaf}.)Qua subi-}
\l{\cel{3}sis kelka perturbo di sua fakultati spiritala, etikala. -}
\l{\cel{3}EFI.}
\l{}
\l{\sou{malakito}. (\sou{mineral.})\cel{2}Lapido verda, ek kupro karbonatoza,}
\l{\cel{3}apta divenar bele polisita, uzata da la juvelisti.- DEFIRS.}
\l{}
\l{\sou{malario}. (\sou{patol.}) Marsho-febro.-}
\l{}
\l{\sou{maldicensar}. (\sou{netrans.}) (\sou{pri)}. Dicar, pri ulu, lo mala quan}
\l{\cel{3}onu savas pri lu. - FIS.}
\l{}
\l{\sou{malear}.\cel{2}(\sou{trans}.) Extensar (ordinare: metalo) per martelagar}
\l{\cel{3}(sen varmigo preparala) aden lami o folii plu o min dina.}
\l{\cel{3}- EFIS.}
\l{}
\l{\sou{meledikar}. (\sou{trans.)}\cel{1}I.Advokar mala chanco kontre ulu.- II.}
\l{\cel{3}(\sou{aludante patrulo}). Advokar la iraco deala (kontre filiulo}
\l{\cel{3}kruelega, senkordia.) - III. (\sou{aludante sacerdoto}). Advokar,}
\l{\cel{3}per vorti, per ritui religiala, la blamo deala kontre ulu,}
\l{\cel{3}kontre ulo. - IV. (\sou{aludante Deo}). Destinar a la blamo}
\l{\cel{3}eterna. - DEFIS.}
\l{}
\l{\sou{maleolo}. (\sou{anat.}) Saliajo osta, ye la infra extremajo di singla}
\l{\cel{3}ek la gambi - interna-latere, forme di salio di la tibio;}
\l{\cel{3}extera-latere, forme di salio di la extremajo di la peroneo.}
\l{\cel{3}- EFIS.}
\l{}
\l{malgre. Sen impedesar da la nevolo quan opozas ulu, da la}
\l{\cel{3}obstaklo quan opozas ulo. - FI.}
\l{}
\l{\sou{malio}. Ludo ye qua onu avancigas bulo ek busligno, per quaza}
\l{\cel{3}martelo ek ligno harda, ferizita, kun mancho longa e flex-}
\l{\cel{3}ebla. - EFIS.}
\l{}
\l{malico. I. Plezuro de agar ulo mala. - II. Su-amuzo de agar}
\l{\cel{3}o dicar malignalaji negrava. - EF.}
}
